\section{Conclusions}
On the 32 additional ChronoGraph services, A3 reduced parameter count from 74 for E5 to 15 and improved memory MASE by 0.02523 under the prespecified fixed-cohort contrast (ten favorable seed means; Holm-adjusted $p=0.0078125$). It did not outperform DLinear in the primary mean, and its CPU mean was worse than both primary references. The small PatchTST configuration had the lowest observed primary mean MASE for CPU and memory, without a prespecified significance claim against NLinear. In the external first-day case, persistence had the lowest MASE for both resources. These results favor careful baseline selection and domain-sensitive interpretation over a universal missingness-architecture claim.

Further research requires independently timed cohorts and longer external periods, documented release preprocessing and graph availability, real interruption measurements, matched component interventions and measured controller outcomes. The local archive permits reconstruction of the numerical comparisons and their limits. Editorial use still requires both authors' final approval, a corresponding-author designation and decisions on prior-publication disclosure and authorized data/code release.
The historical phases supply a necessary part of this answer: E5's paired reduction was positive for memory (1.280\%) and negative for CPU ($-0.984\%$), and V2 did not retain the full memory benefit. The secondary extensions show increasing errors at longer direct horizons and mixed responses to larger histories. These findings make simple references and negative results central to model assessment rather than incidental exceptions.
